% ---- pg-334 （书页 322） ----

首先, 我们建立一个反映天然气输送量 $(Y)$ 与 DHD、风速和周末虚拟变量之间关系的回归模型. 其他可能对天然气消费有某种影响的变量, 如云覆盖百分比等均被吸收到误差项中. 经过几次尝试性拟合后, 我们决定在模型中不仅要加进当前的 DHD 值, 而且加进前一天的 DHD 值(在\hyperref[tab:7.4]{表 7.4} 中, DHD 滞后值表示前一天的 DHD 值). 得到的拟合模型为
\[
\begin{aligned}
\text{天然气输送量}&=1.858+5.874\ \mathrm{DHD}+1.405\ \mathrm{DHD}\ \text{滞后值}\\
&\quad+1.315\ \text{风速}-15.857\ \text{周末}
\end{aligned}
\]
决定系数为 $R^2=0.952$. 除去常数项, 所有的系数均显著. 拟合看起来很好(模型中的常数项可以摒弃, 这对结果不会有实质性影响). 然而, 如果我们计算相邻两期残差之间的相关系数, 即计算一阶自相关系数, 我们却得到
\[
\text{一阶自相关系数}=r_1(\hat{\epsilon})=\frac{\sum_{j=2}^{n}\hat{\epsilon}_j\hat{\epsilon}_{j-1}}{\sum_{j=1}^{n}\hat{\epsilon}_j^2}=0.52
\]
0.52 这个自相关系数值太大, 不可忽视. 前 15 阶残差自相关图表明, 在相隔 7 期(一周)的诸误差之间也许还存在某种相依性. 这种相依性使模型系数的 $t$ 检验和与之相联系的 $P$ 值无效.

修正这个模型的第一步, 是将回归模型中的独立误差用可能有的相依噪音项序列 $N_j$ 代替. 也就是说, 我们对 $N_j$ 构造一个模型, 将每一个 $N_j$ 与其前一天的值 $N_{j-1}$, 前一周的值 $N_{j-7}$, 以及独立误差 $\epsilon_j$ 建立联系. 这样, 我们考虑
\[
N_j=\phi_1N_{j-1}+\phi_7N_{j-7}+\epsilon_j
\]
其中 $\epsilon_j$ 为独立正态随机变量, 均值为 0, 方差为 $\sigma_\epsilon^2$. $N_j$ 所满足上述的方程形式称为自回归模型(见文献[8]). 对天然气输送量的这一组合式回归模型(其噪音部分采用自回归模型), 其计算的 SAS 命令和部分拟合结果的输出见显示板 7.3.

拟合的模型为
\[
\begin{aligned}
\text{天然气输出}&=2.130+5.810\ \mathrm{DHD}+1.426\ \mathrm{DHD}\ \text{滞后值}\\
&\quad+1.207\ \text{风速}-10.109\ \text{周末}
\end{aligned}
\]
而噪音项的时间相依性由下式估计:
\[
N_j=0.470N_{j-1}+0.240N_{j-7}+\epsilon_j
\]
$\epsilon$ 的方差估计值为 $\hat{\sigma}_\epsilon^2=228.89$.

由显示板 7.3 我们看到, 在扩充的模型中, 残差间的自相关系数全部可以忽略. 每个残差自相关系数都落在位于 0 两侧的两个标准误差估计值之间. 还有, 一组阶数接连增大的残差自相关系数的加权平方和并不大, 这可根据这个统计量的 $P$ 值来判断. 也就是说, 没有理由拒绝假设“一组阶数接连增大的自相关系数同时为 0”. 在显示板 7.3 中, 考察了自相关阶数为 $1\sim6, 1\sim12, 1\sim18$ 和 $1\sim24$ 的四个组.

现已对噪音恰当地建立了模型, 对每个预测量的系数以及回归的显著性等等的检验, 现在都成立.\footnote{{\footnotesize 这些检验是用额外平方和方法得出的, 不过应将这种方法应用于回归加自回归噪音模型. 这就是计算机输出中所描述的检验.}} 常数项在最终模型中可以略去. 当这样做时, 在最终模型中几乎不会引起什么变
